\chapter[SYMBOLES APRS]{Symboles APRS}\label{symboles-aprs}

\section{Trois méthodes}\label{trois-muxe9thodes}

Pour spécifier un symbole (icône d'affichage) : - Dans le champ
Information AX.25 (préféré). - Dans l'adresse de destination AX.25. -
Dans le SSID de l'adresse source AX.25.

Les deux dernières pour les trackers autonomes qui ne peuvent pas mettre
le symbole dans Information.

\section{Les tables}\label{les-tables}

Deux tables : \textbf{Primary} et \textbf{Alternate}. Voir annexe 2 pour
le listing complet.

Différence : certains symboles de la table alternative peuvent être
overlaidés d'un caractère alphanumérique (marqués « {[}with overlay{]}
»). Aucun symbole de la table primaire ne peut l'être.

Chaque symbole codé de 3 manières : - \texttt{/\$} ou
\texttt{\textbackslash{}\$} --- Information field. - \texttt{GPSxyz} ---
destination générique. - \texttt{GPSCnn} ou \texttt{GPSEnn} --- autre
forme de destination générique.

15 symboles de la table primaire ont un SSID associé (ex : petit avion =
\texttt{-7}), pour les trackers autonomes.

\section{Symboles dans le champ
Information}\label{symboles-dans-le-champ-information}

Symbole = Symbol Table Identifier (1 char) + Symbol Code (1 char).

Ex : \texttt{@092345z4903.50N/07201.75W\textgreater{}088/036…} ---
\texttt{/} = Sym Table ID (primary), \texttt{\textgreater{}} = Symbol
Code (voiture).

\begin{longtable}[]{@{}
  >{\raggedright\arraybackslash}p{(\columnwidth - 2\tabcolsep) * \real{0.5000}}
  >{\raggedright\arraybackslash}p{(\columnwidth - 2\tabcolsep) * \real{0.5000}}@{}}
\toprule\noalign{}
\begin{minipage}[b]{\linewidth}\raggedright
ID
\end{minipage} & \begin{minipage}[b]{\linewidth}\raggedright
Signification
\end{minipage} \\
\midrule\noalign{}
\endhead
\bottomrule\noalign{}
\endlastfoot
\texttt{/} & Table primaire (surtout stations) \\
\texttt{\textbackslash{}} & Table alternative (surtout objets) \\
\texttt{0-9} & Overlay numérique, table alternative (uncompressed
lat/long) \\
\texttt{a-j} & Overlay numérique, table alternative (compressed
lat/long) --- \texttt{a-j} → \texttt{0-9} \\
\texttt{A-Z} & Overlay alpha, table alternative \\
\end{longtable}

\section{Overlays dans le champ
Information}\label{overlays-dans-le-champ-information}

Si Sym Table ID = \texttt{0-9} ou \texttt{A-Z} (ou \texttt{a-j} en
compressed), le symbole vient de la table alternative avec le caractère
d'overlay.

Ex uncompressed : \texttt{@092345z4903.50N307201.75W\textgreater{}…} ---
chiffre \texttt{3} overlaidé sur icône voiture.

Ex compressed avec lettre :
\texttt{=BL!!\textless{}*e7\ \textgreater{}7P{[}} --- voiture overlaidée
\texttt{B}.

\textbf{Attention} : en compressed, Sym Table ID numérique interdit (les
positions compressées ne commencent jamais par un chiffre). Utiliser
\texttt{a-j} mappée à \texttt{0-9}.

Ex : \texttt{=d5L!!\textless{}*e7\ \textgreater{}7P{[}} --- \texttt{d}
mappe à overlay \texttt{3}.

Un symbole overlay-capable n'est pas obligé d'avoir un overlay.

\section{Symboles dans l'adresse de
destination}\label{symboles-dans-ladresse-de-destination}

Destinations génériques : {\ttfamily \lit{GPS}xyz}, {\ttfamily \lit{GPS}Cnn},
{\ttfamily \lit{GPS}Enn}, {\ttfamily \lit{SPC}xyz}, {\ttfamily \lit{SYM}xyz}. - \texttt{xy} /
\texttt{nn} = entrées des tables. - \texttt{z} = overlay (si permis) ou
espace filler.

{\ttfamily \lit{GPS}/\lit{SPC}/\lit{SYM}xy} et {\ttfamily \lit{GPS}Cnn/\lit{GPS}Enn} interchangeables. Ex :
{\ttfamily \lit{GPS}BM}, {\ttfamily \lit{SPC}BM}, {\ttfamily \lit{SYM}BM}, {\ttfamily \lit{GPS}C12}
désignent tous « Boy Scouts » (primary).

\section{Overlays dans la
destination}\label{overlays-dans-la-destination}

Si \texttt{z} non espace dans {\ttfamily \lit{GPS}xyz} = overlay \texttt{0-9} ou
\texttt{A-Z}. Seuls les symboles alternate marqués « {[}with overlay{]}
». Ex : voiture overlaidée \texttt{3} → {\ttfamily \lit{GPS}N\ 3}.
{\ttfamily \lit{GPS}Cnn} et {\ttfamily \lit{GPS}Enn} ne peuvent pas être overlaidés.

\section{Symbole dans le SSID
Source}\label{symbole-dans-le-ssid-source}

\begin{longtable}[]{@{}llll@{}}
\toprule\noalign{}
SSID & Icône & SSID & Icône \\
\midrule\noalign{}
\endhead
\bottomrule\noalign{}
\endlastfoot
\texttt{-0} & {[}pas d'icône{]} & \texttt{-8} & Ship (power boat) \\
\texttt{-1} & Ambulance & \texttt{-9} & Car \\
\texttt{-2} & Bus & \texttt{-10} & Motorcycle \\
\texttt{-3} & Fire Truck & \texttt{-11} & Balloon \\
\texttt{-4} & Bicycle & \texttt{-12} & Jeep \\
\texttt{-5} & Yacht & \texttt{-13} & Recreational Vehicle \\
\texttt{-6} & Helicopter & \texttt{-14} & Truck \\
\texttt{-7} & Small Aircraft & \texttt{-15} & Van \\
\end{longtable}

\section{Précédence des
symboles}\label{pruxe9cuxe9dence-des-symboles}

Un paquet APRS ne devrait pas contenir plus d'un symbole. Si erreur,
précédence :

\begin{enumerate}
\def\labelenumi{\arabic{enumi}.}
\tightlist
\item
  Symbole dans Information field (le plus prioritaire).
\item
  Sinon symbole dans destination.
\item
  Sinon SSID source.
\end{enumerate}

\begin{center}\rule{0.5\linewidth}{0.5pt}\end{center}
